\documentclass{beamer}

\usetheme{Berlin}
\usecolortheme{Orchid}
\useoutertheme{miniframes}
\setbeamertemplate{navigation symbols}{}

\usepackage[utf8]{inputenc}
\usepackage[T1]{fontenc}
\usepackage{amsmath}
\usepackage{amssymb}
\usepackage{booktabs}
\usepackage{graphicx}
\graphicspath{{../images/}}

\usepackage{xcolor}
\usepackage{listings}
\usepackage{hyperref}
\usepackage{tikz}
\usetikzlibrary{arrows.meta,positioning,shapes.geometric,calc}

\lstset{
  language=Java,
  basicstyle=\ttfamily\small,
  keywordstyle=\color{blue},
  commentstyle=\color{gray},
  stringstyle=\color{teal},
  showstringspaces=false,
  columns=fullflexible,
  breaklines=true
}

% Per-slide source line (references shown on the slide itself).
\newcommand{\slidesources}[1]{%
  \vfill
  {\tiny\color{gray}\raggedright\sloppy\parbox{\linewidth}{Sources: #1}\par}%
}

% Unit-wise source strings for this subject (used by slides/notes).
% Path is relative to the lecture `latex/` folder (the compilation CWD).
% OOPS with Java (CSEG1044) — unit-wise sources
%
% These are short, student-facing reference strings intended to be used with:
%   - \slidesources{...}  (in beamer slides)
%   - \notesources{...}   (in lecture notes)
%
% Style inspiration: other subjects in My-Subjects/ use semicolon-separated sources
% and include an "accessed" date for web documentation.

\newcommand{\OOPSUnitISources}{Schildt, \textit{Java: A Beginner's Guide} (10e), McGraw-Hill (Java fundamentals + OOP basics).; Oracle Java Tutorials: Getting Started + Language Basics + Classes and Objects (accessed \today).; Java SE API docs (e.g., \texttt{String}, \texttt{Scanner}) (accessed \today).}

\newcommand{\OOPSUnitIISources}{Schildt, \textit{Java: The Complete Reference} (12e), McGraw-Hill (classes, inheritance, packages, interfaces).; Bloch, \textit{Effective Java} (3e), Addison-Wesley, 2018 (design choices; interfaces vs abstract classes).; Oracle Java Tutorials: Classes and Objects + Interfaces + Packages (accessed \today).; Java SE API docs (e.g., \texttt{Comparable}, \texttt{Runnable}) (accessed \today).}

\newcommand{\OOPSUnitIIISources}{Schildt, \textit{Java: The Complete Reference} (12e) (nested/inner classes, exceptions, threads, I/O).; Oracle Java Tutorials: Nested Classes + Exceptions + Concurrency + Basic I/O (accessed \today).; Java SE API docs (e.g., \texttt{Thread}, \texttt{AutoCloseable}) (accessed \today).}

\newcommand{\OOPSUnitIVSources}{Schildt, \textit{Java: The Complete Reference} (12e) (generics, lambdas, Swing, JDBC).; Bloch, \textit{Effective Java} (3e) (generics + lambdas).; Oracle Java Tutorials: Generics + Lambda Expressions + Swing + JDBC Basics (accessed \today).; Java SE API docs (e.g., \texttt{java.util.function}, \texttt{javax.swing}, \texttt{java.sql}) (accessed \today).}

\newcommand{\OOPSUnitVSources}{Schildt, \textit{Java: The Complete Reference} (12e) (Collections Framework + wrapper classes).; Oracle Java docs: Collections Framework overview (accessed \today).; Java SE API docs (\texttt{java.util} collections; wrapper classes in \texttt{java.lang}) (accessed \today).}

\newcommand{\OOPSUnitVISources}{Git SCM: \textit{Pro Git} (2e) (version control workflow), accessed \today.; GitHub Docs (repositories, issues, pull requests), accessed \today.; Oracle Java Tutorials: Swing + JDBC (accessed \today).; JUnit 5 User Guide (accessed \today).; Apache Maven Project: Getting Started (accessed \today).; Gradle User Manual: Getting Started (accessed \today).; UML class diagrams: UML 2.x overview/spec (accessed \today).}

\newcommand{\OOPSMiscWhyInterfacesSources}{Schildt, \textit{Java: The Complete Reference} (12e) (interfaces).; Bloch, \textit{Effective Java} (3e) (prefer interfaces where appropriate).; Oracle Java Tutorials: Interfaces (accessed \today).; Java SE API docs (e.g., \texttt{Comparable}, \texttt{Runnable}) (accessed \today).}



\title[OOPS with Java]{Object Oriented Programming (CSEG1044)}
\subtitle{Unit 03 -- Lecture 03: Multithreading + Synchronization}
\author{Tofik Ali}
\institute{School of Computer Science, UPES Dehradun}
\date{\today}

\begin{document}

\begin{frame}[fragile]
  \titlepage
  \vspace{0.5em}
  \slidesources{\OOPSUnitIIISources}
\end{frame}

\begin{frame}{Quick Links}
  \centering
  \hyperlink{sec:overview}{\beamerbutton{Overview}}\hspace{0.6em}
  \hyperlink{sec:concepts}{\beamerbutton{Key Topics}}\hspace{0.6em}
  \hyperlink{sec:demo}{\beamerbutton{Demo}}\hspace{0.6em}
  \hyperlink{sec:summary}{\beamerbutton{Summary}}
  \slidesources{\OOPSUnitIIISources}
\end{frame}

\begin{frame}{Agenda}
  \tableofcontents
  \slidesources{\OOPSUnitIIISources}
\end{frame}

\section{Overview}
\label{sec:overview}

\begin{frame}{Learning Outcomes}
  \begin{itemize}[<+->]
    \item Create threads using \texttt{Thread} and \texttt{Runnable}.
    \item Describe thread lifecycle and understand priority limitations.
    \item Identify race conditions and fix them with \texttt{synchronized}.
  \end{itemize}
  \slidesources{\OOPSUnitIIISources}
\end{frame}

\section{Key Topics}
\label{sec:concepts}

\begin{frame}{Unit 03: Nested Classes, Exceptions, Multithreading \& IO Streams}
  \begin{itemize}[<+->]
    \item Lecture focus: Multithreading + Synchronization
  \end{itemize}
  \slidesources{\OOPSUnitIIISources}
\end{frame}

\begin{frame}{Today: Roadmap}
  \begin{itemize}[<+->]
    \item Concurrency motivation; `Thread` class and `Runnable` interface
    \item Thread lifecycle and priorities (what they mean / limitations)
    \item Race conditions and synchronization (`synchronized`, basic patterns)
  \end{itemize}
  \slidesources{\OOPSUnitIIISources}
\end{frame}

\begin{frame}{Concurrency vs Parallelism}
  \begin{itemize}[<+->]
    \item \textbf{Concurrency}: tasks overlap in time (good for responsiveness).
    \item \textbf{Parallelism}: tasks run simultaneously on multiple cores (good for speed).
  \end{itemize}
  \slidesources{\OOPSUnitIIISources}
\end{frame}

\begin{frame}{Creating threads (two ways)}
  \begin{itemize}[<+->]
    \item Extend \texttt{Thread} and override \texttt{run()}.
    \item Implement \texttt{Runnable} and pass it to \texttt{new Thread(r)}.
    \item \textbf{Important:} call \texttt{start()} to create a new thread.
  \end{itemize}
  \slidesources{\OOPSUnitIIISources}
\end{frame}

\begin{frame}[fragile]{\texttt{start()} vs \texttt{run()} (very important)}
\begin{lstlisting}
Runnable task = new Runnable() {
    public void run() {
        System.out.println(Thread.currentThread().getName());
    }
};

Thread t = new Thread(task, "worker");

t.run();   // runs on current thread (main)
t.start(); // runs on a new thread (worker)
\end{lstlisting}
  \begin{itemize}[<+->]
    \item \texttt{run()} is a normal method call (no new thread).
    \item \texttt{start()} asks JVM/OS to schedule a new thread.
  \end{itemize}
  \slidesources{\OOPSUnitIIISources}
\end{frame}

\begin{frame}[fragile]{Waiting for threads: \texttt{join()}}
\begin{lstlisting}
Runnable task = new Runnable() {
    public void run() { /* do work */ }
};

Thread t1 = new Thread(task);
Thread t2 = new Thread(task);
t1.start(); t2.start();

t1.join();  // wait for t1 to finish
t2.join();  // wait for t2 to finish
System.out.println("Both done");
\end{lstlisting}
  \begin{itemize}[<+->]
    \item Without \texttt{join()}, main may finish early.
    \item \texttt{join()} throws \texttt{InterruptedException}.
  \end{itemize}
  \slidesources{\OOPSUnitIIISources}
\end{frame}

\begin{frame}{Lifecycle + priorities (practical)}
  \begin{itemize}[<+->]
    \item New $\rightarrow$ Runnable/Running $\rightarrow$ Waiting/Blocked $\rightarrow$ Terminated
    \item Waiting happens due to sleep, locks, I/O, etc.
    \item Priority is a \textbf{hint}, not a guarantee (depends on OS/JVM).
  \end{itemize}
  \slidesources{\OOPSUnitIIISources}
\end{frame}

\begin{frame}{Race condition (core problem)}
  \begin{itemize}[<+->]
    \item Multiple threads modify shared state at the same time.
    \item Result depends on timing/interleaving.
    \item Example: \texttt{value++} is not atomic (read-modify-write).
  \end{itemize}
  \slidesources{\OOPSUnitIIISources}
\end{frame}

\begin{frame}{Synchronization with \texttt{synchronized}}
  \begin{itemize}[<+->]
    \item Protect a critical section using a lock/monitor.
    \item Only one thread at a time can execute synchronized code on the same lock.
    \item Keep synchronized sections small; avoid deadlocks.
  \end{itemize}
  \slidesources{\OOPSUnitIIISources}
\end{frame}


\section{Demo / Example}
\label{sec:demo}

\begin{frame}[fragile]{Example: synchronized counter}
\begin{lstlisting}
class SafeCounter {
    private int value = 0;
    synchronized void inc() { value++; }
    synchronized int get() { return value; }
}
\end{lstlisting}
  \slidesources{\OOPSUnitIIISources}
\end{frame}

\begin{frame}[fragile]{Thread demo (two threads) (1/2)}
\begin{lstlisting}
public class Demo {
    public static void main(String[] args) throws InterruptedException {
        SafeCounter c = new SafeCounter();

        Runnable task = new Runnable() {
            public void run() {
                for (int i = 0; i < 100000; i++) c.inc();
            }
        };
\end{lstlisting}
  \slidesources{\OOPSUnitIIISources}
\end{frame}

\begin{frame}[fragile]{Thread demo (two threads) (2/2)}
\begin{lstlisting}
// continued...
        Thread t1 = new Thread(task);
        Thread t2 = new Thread(task);
        t1.start();
        t2.start();
        t1.join(); t2.join();
        System.out.println(c.get()); // expected 200000
    }
}
\end{lstlisting}
  \slidesources{\OOPSUnitIIISources}
\end{frame}

\begin{frame}{Mini Exercises (2--3 mins)}
  \begin{enumerate}[<+->]
    \item Difference: \texttt{start()} vs \texttt{run()}?
    \item Why is \texttt{value++} unsafe in multiple threads?
    \item What does \texttt{synchronized} guarantee (in one sentence)?
  \end{enumerate}
  \slidesources{\OOPSUnitIIISources}
\end{frame}


\section{Summary}
\label{sec:summary}

\begin{frame}{Key Takeaways}
  \begin{itemize}[<+->]
    \item Threads enable concurrency for responsiveness and throughput.
    \item Race conditions occur with shared mutable state.
    \item \texttt{synchronized} protects critical sections using a lock.
  \end{itemize}
  \slidesources{\OOPSUnitIIISources}
\end{frame}

\begin{frame}{Checkpoint Questions}
  \begin{enumerate}[<+->]
    \item What is a race condition?
    \item Why is thread priority not a guarantee?
    \item How does \texttt{synchronized} help?
  \end{enumerate}
  \slidesources{\OOPSUnitIIISources}
\end{frame}

\begin{frame}{Next Lecture Preview}
  \textbf{Next:} Java I/O Streams + FileReader/FileWriter
  \vspace{0.6em}
  \begin{itemize}[<+->]
    \item We will read/write files and handle errors safely with exceptions.
  \end{itemize}
  \slidesources{\OOPSUnitIIISources}
\end{frame}

\end{document}
